\section*{Course Logistics}
This is a 4-credit course requiring CSSE 120 as a prerequisite. 
See the \href{https://catalog.rose-hulman.edu/course-search/?details&srcdb=2025&code=CSSE%20220}{course catalog} 
for more details. 
% this is a 4-credit course with a prerequisite of 
% CSSE 120 (or demonstrated programming experience). 
% The course follows a ``$3 \times 2$'' combined lecture/lab format. %, 
% meeting 3 days per week, 2 periods each day. 
% Each day will have practice exercises, often with group discussions. 
% This in-class practice is extremely valuable for learning, 
% so regular attendance is expected. 

\subsection*{Time{(s)} and Location{(s)}:} 
Both sections meet MWF in Olin 157 (O157). 
Section 01 is 8:00--9:50 a.m., and 
Section 04 is 3:00--4:50 p.m. 

\subsection*{Course Description:} 
Object-oriented programming concepts, 
including the use of inheritance, interfaces, polymorphism, abstract data types,
and encapsulation to enable software reuse and assist in software maintenance. 
Recursion, GUIs and event handing. 
Use of common object-based data structures, 
including stacks, queues, lists, trees, sets, maps, and hash tables. 
Space/time efficiency analysis. Testing. Introduction to {UML}. 

\subsection*{Learning Outcomes:} 
Students who successfully complete this course should be able to:
\begin{enumerate}[label=\arabic*.,noitemsep,nolistsep]
\item Develop software that incorporates the following techniques:
    \begin{enumerate}[label={(\alph*)},noitemsep,nolistsep]
    \item Inheritance and class hierarchies
    \item Interfaces
    \item Polymorphism
    \item Casting
    \item Exceptions
    \item Collections
    \item Event-driven graphical user interfaces
    \item Exploring and using large-scale API packages (e.g. GUI libraries)
    \item Recursion
    \end{enumerate}
\item Perform the following steps of the software development cycle effectively:
    \begin{enumerate}[label={(\alph*)},noitemsep,nolistsep]
    \item Design expressed as UML class diagrams
    \item Documentation before coding
    \item Unit and system testing
    \end{enumerate}
\item Explain the implementation of linked lists
\item Analyze the asymptotic worst case run time of simple programs.
\item Predict the performance of simple algorithms, including search and sort, given their asymptotic worst, best, and average case run times.
\item Select basic data structures (i.e., arrays, linked lists) based on asymptotic time complexity of typical operations.
\item Work in a team of 2--4 students on a small-to-medium-size software development project including at least three iterative development cycles, demonstrating effective:
    \begin{enumerate}[label={(\alph*)},noitemsep,nolistsep]
    \item Use of team roles
    \item Team decision making
    \item Division of labor
    \end{enumerate}
\item Identify opportunities and limitations of using 
    AI tools in the software development cycle. 
\end{enumerate}

